\begin{frame}
	\frametitle{O neurônio \textit{Leaky Integrate-and-Fire} (LIF)}
	\only<1>{
		\framesubtitle{O modelo RC}
		\begin{columns}
			\column{.48\linewidth}
			\begin{itemize}
				\item O LIF se comporta como um circuito Resistor-Capacitor (RC).
				\item $R$: resistência de vazamento; $C$: capacitância; $V_{mem}$: potencial acumulado; $\vartheta$: interruptor que descarrega o capacitor (dispara um pulso) ao atingir o limiar.
				\item Governado por $\tau\frac{dV_{mem}}{dt} = R\,I_{in}(t) - V_{mem}(t)$, com $\tau = R\cdot C$ \cite{eshraghian2023training}.
			\end{itemize}
			\column{.48\linewidth}
			\begin{center}
				\includegraphics[width=.85\linewidth]{../../00-thesis/monography/images/rcmodel}\par
				{\tiny Fonte: \cite{eshraghian2023training}}
			\end{center}
		\end{columns}
	}
	\only<2>{
		\framesubtitle{Exemplo numérico: decaimento sem entrada}
		\par Com $R=5\,\Omega$, $C=1\,F$ (logo $\tau = R\cdot C = 5$) e potencial inicial $V_{mem}(0)=1$, sem corrente de entrada:
		\begin{equation*}
			V_{mem}(t) = V_{mem}(0)\cdot e^{-t/\tau}
			\quad\Longrightarrow\quad
			V_{mem}(1) = e^{-1/5} \approx 0{,}8187
		\end{equation*}
		\par O potencial \textbf{vaza} exponencialmente na ausência de estímulo --- daí o nome \textit{leaky}. Quando um estímulo chega, o potencial sobe (mesma matemática, sinal oposto) até cruzar um limiar $V_{th}$, disparando um pulso e reiniciando.
	}
	\only<3>{
		\framesubtitle{Comportamento completo simulado}
		\begin{center}
			\includegraphics[width=.62\linewidth]{../../00-thesis/monography/images/membranePotentialFull}\par
			{\tiny Corrente de 0,5 mA injetada entre os passos 51 e 70. Fonte: elaborado pelo autor.}
		\end{center}
	}
\end{frame}
